

\section{Extensions}
Our framework naturally allows for some extensions and future directions, of which we here outline two. As mentioned before, our framework in this paper only uses the binary value of $\Phi$ -- it is either zero or non-zero -- but the true strength of IIT lies in its \emph{size}. Hence our first extension constructs a proper analogue to the $\Phi$-value in our framework, which we denote $\Phi_{coh}$. Further, our cohomological $\Phi$-value, as well as our previously defined obstruction classes, all lie in the first cohomology groups $H^1_{int}(S)$ of the system $(S, \D, T)$. Utilizing the higher order cohomology groups as well, we obtain higher order obstruction classes and higher order cohomological $\Phi$-values, which are analogous to the higher order synergy anf contextuality of \citeme. 

\subsection{Cohomological integrated information}
In \cref{ssec:obstruction-and-IIT} we introduced a distance function between states, 
\[d\: \St(S)\times \St(S) \to \R\geqz,\]
that we used to describe Tull--Kleiner's genealized IIT-algorithm. We here use this to infer a distance function on the cohomology groups, which we then use to define te cohomological integrated information $\Phi_{coh}$. 

\begin{definition}
    Let $A\in \Sub_s(S)$ be a subsystem. The \emph{distance} between two endomorphisms $f,g$ of $A$ is defined by how differently they act on the local state $s_{\mid A}$. In other words, 
    \[\Vert f-g \Vert_A := d(f\circ s_{\mid A}, g\circ s_{\mid A}).\]
    This extends linearly to a distance function on $\F(A)$. Furthermore, given a $k$-cochain $c \in C^k(S;\F)$, we define its \emph{size} as its largest value on the set of all $k$-simplicies: 
    \[\Vert c \Vert := \max_\sigma \Vert c(\sigma)\Vert.\]
\end{definition}

In order to infer a size on cohomology classes, we simply need to incorporate the fact that cohomology classes are only defined up to a coboundary, as they are quotient classes of abelian groups. Hence, given a cohomology class $[c]\in H^k_{int}(S)$, we define 
\[\Vert [c] \Vert := \min_{\theta\in C^{k-1}(S;\F)} \Vert c + \delta \theta \Vert.\]
This definition makes the value intrinsic, as it minimizes over all possible representations of the class. 

\begin{definition}
    Let $S$ be a system in a state $s$. The \emph{cohomological integrated information} is defined as the size of the smalles obstruction class, 
    \[\Phi_{coh}(S,s) := \min_{\psi, (\alpha_k)} \Vert [\Delta^{\psi, (\alpha_k)}]\Vert .\]
\end{definition}

The following result immediatly follows from the definition combined with \cref{thm:obstructions-and-IIT}. 

\begin{proposition}
    If $s$ is a state of a system $(S, \D, T)$ such that $\Phi(S,s) = 0$, then also $\Phi_{coh}(S,s)=0$.  
\end{proposition}

We hope to revisit this definition in future work and connect it to the definition of the standard integrated information value $\Phi(S,s)$. For now we explain how this solves Tull and Kleiners search for a functorial construction of integrated information. 

In order to have functoriality, we need to understand what categories are present. In particular, we need a category that describes both systems $(S, \D, T)$ and their states $s$. The recipient category will be the poset category $\R\geqz$ of non-negative real numbers. A map into this poset category is functorial when it is monotone. 

\begin{definition}
    A \emph{system state} in a process theory $\C$, is a tuple $(S,\D, T, s)$, where $(S, \D, T)$ is a system in $\C$ and $s\: \1 \to S$ is a state of $S$. A \emph{morphism} of system states is a causal map 
    \[(S, \D, T, s) \to (S', \D', T', s')\]
    that
    \begin{enumerate}
        \item preserves states: $f \circ s = s'$;
        \item preserves dynamics: $f \circ T = T'$;
        \item preserves decompositions: for a decomposition $(A,B,\psi)$ these exists a decomposition $(f(A), f(B), \psi')$ into subsystems identified by $f$;
        \item induces a functor on the poset categories: $f_* \: \Sub_s(S) \to \Sub_{s'}(S')$;
        \item is compatible with the restriction maps: $f(\rho_U^V(c)) = \rho_{f(U)}^{f(V)(f(c))}$;
        \item is non-expansive: $\Vert f_*(c) \Vert \leq \Vert c \Vert$.
    \end{enumerate} 
    The resulting category of system states and morphisms is denoted $\Sys(\C)$. 
\end{definition}

It follows from the definition of the category above that the cohomological integrated information is a functor, as the obstruction classes are functorial by design, and the properties of morphisms in $\Sys(\C)$ ensures that it is a monotone map. Hence we have the following proposition. 

\begin{proposition}
    \label{prop:functoriality}
    Let $\C$ be a process theory. The cohomological integrated information is a functor 
    \[\Phi_{coh} \: \Sys(\C) \to \R\geqz.\]
    In particular, isomorphic system states have equal cohomological integrated information, 
\end{proposition}

\subsection{Higher order integration}

Our setup of using obstruction theory easily allows us to extend integrated information to higher orders -- something that cohomology theory is a natural arena for. There are some attempts at this already, se for example \citeme, but we will save comparisons for future work. 

Higher obstruction classes are defined using more refined decomposition of the system $S$. The intuition is that, given a decomposition $\psi\: S\simeq A\otimes B\otimes C$, the $\Phi$-values computed pairwise from these could all be zero, even if the total $\Phi$-value is not. This failure should be measured not by the bipartite obstruction classes previously defined, but by a higher obstruction class. 

\begin{definition}
    Let $S$ be an object in a process theory $\C$. A \emph{$k$-ary decomposition} of $S$ is a $k$-tuple $(A_1, A_2, \ldots, A_k)$ together with a causal isomorphism 
    \[\psi \: S \simeq A_1\otimes A_2 \otimes \dots \otimes A_k.\] 
    We will denote such a decomposition simply by $(A_k)$. In particular, a $k$-ary decomposition is an iterated binary composition, hence naturally lives in $\D(S)$. 
\end{definition}

By inserting noise, we can use these to construct \emph{multi-cut} time evolutions $T^{\psi, (\alpha_k)}$, where $(\alpha_k)$ is a sequence of $k$ directions, also allowing for some of the values being empty.  

\begin{definition}
    Given a system $(S, \D, T)$, a $k$-ary decomposition $(A_k)$ and a cut-vector $(\alpha_k)$, we define the associated obstruction class on a $k+1$-simplex $\sigma = U_0\rightarrow \dots \rightarrow U_k$ as the cohomology class of
    \[\Delta^{\psi, (\alpha_k)}(\sigma) = \rho_{U_0}^{U_k}(T_{\mid U_k}) - \rho_{U_0}^{U_k}(T^{\psi, (\alpha_k)}_{\mid U_k}),\]
    where $\rho_{U_0}^{U_k}$ is the iterated restriction from $U_k$ to $U_0$ defined by $\sigma$. 
\end{definition}

This is a class living in $H^k_{int}(S)$, which measures the failure of the largest subsystem $U_k$ to decompose into $k$ parts that are ``$(\alpha_k)$-independent'' -- if the direction vector consisted of only bidirectional arrows, then these parts would have been independent. For the case $k=1$, these classes reduce to the ones studied earlier. 

Analogously to the construction above for $\Phi_{coh}(S,s)$ we can define higher order cohomological integrated information. 

\begin{definition}
    Let $(S,\D,T,s)$ be a system state. Its \emph{$k$-th order cohomological integrated information} is defined as the size of the smallest higher order obstruction class: 
    \[\Phi_{coh}^k(S,s) := \min_{\psi, (\alpha_k)}\Vert [\Delta^{\psi, (\alpha_k)}]\Vert.\]
\end{definition}

As in \cref{prop:functoriality}, these higher order classes are functorial, and for $k=1$ we have $\Phi^1_{coh}(S,s)=\Phi_{coh}(S,s)$ as defined above. 

This defines for a system state $(S,\D,T,s)$ a sequence of real numbers $\Phi^1_{coh}(S,s), \Phi^2_{coh}(S,s), \Phi^3_{coh}(S,s), \ldots$ that detect higher and higher orders of integration. We round off the paper with the following defintion and conjecture.  

\begin{definition}
    The integration order of a system state $(S,\D,T,s)$ is defined as the smallest $k$ such that $\Phi_{coh}^k(S,s)\neq 0$. If all higher cohomological integrated information vanishes, i.e., $\Phi_{coh}^k(S,s)=0$ for all $k$, then we say it has integration order $\infty$.  
\end{definition}

\begin{conjecture}
    If $(S,\D,T,s)$ is a system state with integration order $\infty$, then $\Phi(S,s)=0$. 
\end{conjecture}

